% Generated draft: requires mathematical and source review.
\documentclass[11pt]{article}
\usepackage[margin=25mm]{geometry}
\usepackage{fontspec}
\setmainfont{Latin Modern Roman}
\usepackage{amsmath,amssymb,mathtools,unicode-math}
\setmathfont{Latin Modern Math}
\usepackage{enumitem,needspace,xurl,hyperref}
\setlength{\parindent}{0pt}
\setlength{\parskip}{5pt}
\setlength{\emergencystretch}{3em}
\newcommand{\N}{\mathbb{N}}
\newcommand{\Z}{\mathbb{Z}}
\newcommand{\Q}{\mathbb{Q}}
\newcommand{\R}{\mathbb{R}}
\newcommand{\C}{\mathbb{C}}
\newcommand{\F}{\mathbb{F}}
\newcommand{\A}{\mathbb{A}}
\newcommand{\PP}{\mathbb{P}}
\newcommand{\E}{\mathbb{E}}
\newcommand{\eps}{\varepsilon}
\newcommand{\dd}{\,\mathrm{d}}
\begin{document}
\section*{UnsolvedMath ID: 30001695}
\textbf{Source problem code:} catalogue entry
\subsection{Definitions and mathematical statement}
Use the conventional finite square-grid definition of a polyomino: let \(S\subset\Z^2\) be finite and nonempty, with any two elements joined by a sequence of elements of \(S\) whose successive differences belong to \(\{(1,0),(-1,0),(0,1),(0,-1)\}\). Set
\[
P=\bigcup_{(m,n)\in S}[m,m+1]\times[n,n+1],\qquad J=\partial P.
\]
Require that \(J\) be a Jordan curve, meaning that it is the image of a continuous injective map from the unit circle into \(\R^2\).

The proposed assertion is that for every such \(P\), there exist \(x,u,v\in\R^2\) satisfying
\[
\|u\|=\|v\|>0,\qquad u\cdot v=0,\qquad
\{x,x+u,x+u+v,x+v\}\subseteq J.
\]
Thus the four distinct points are the vertices of a nondegenerate Euclidean square. They may lie anywhere on \(J\); no lattice-point, boundary-corner, axis-alignment, or square-interior containment condition is imposed by the supplied statement.

The record does not explicitly define ``polyomino.'' The conventional definition above is an interpretive convention requiring confirmation against the original report, particularly if a more restrictive grid variant was intended.
\subsection{Short English statement}
Does every polyomino whose entire boundary is a Jordan curve have four boundary points forming a nondegenerate square? The draft uses the conventional finite, edge-connected square-grid meaning of polyomino; the precise intended grid conventions require source review.
\subsection{Sources}
\begin{itemize}
\item Local UnsolvedMath registry, supplied record with permanent ID \(30001695\), source problem code \texttt{catalogue entry}, title \emph{Square-Peg Problem for Polyomino Boundaries}; collection: Oberwolfach Reports Open Problems. The supplied proposed year is 2011; no individual proposer is identified.
\item Supplied source citation and literature reference labelled ``original OWR report'': \emph{Topological and Geometric Combinatorics} (2011), Oberwolfach Reports, pp.\ 349--423, DOI 10.4171/owr/2011/08. \url{https://doi.org/10.4171/owr/2011/08}. The supplied background describes these reports as documenting workshops and their open-problem discussions.
\item Supplied literature reference labelled ``current square-peg progress'': \url{https://arxiv.org/abs/2301.01340}. No author, title, or specific supporting passage is supplied for this reference.
\item Literature assessment checked at \texttt{2026-08-21}: the supplied curation labels the problem open, states that the underlying Toeplitz square-peg problem remains open for arbitrary Jordan curves, and reports that the grid/polyomino variant was not resolved by the located literature. This is supplied dated curation derived from the OWR record, not an independently verified assessment; time-sensitive status and cited claims require independent checking.
\item The registry labels the statement \texttt{corrected\_verified}, with a supplied claim of manual semantic, typographic, and LaTeX review confirming a meaning-preserving, self-contained cleanup. This editorial label is recorded as provenance and does not constitute independent mathematical verification of this draft.
\end{itemize}
\end{document}
